\documentclass[12pt]{article}
\usepackage[utf8]{inputenc}
\usepackage[T1]{fontenc}
\usepackage{hyperref}
\title{Yil Language Resources for K-12 Curriculum Development}
\author{Compiled by Claude Sonnet 5.5 (claude-sonnet-5-5) for Andr\'as Kornai}
\date{}
\begin{document}
\maketitle

\section{Introduction}
Yil (ISO 639-3: yll) is a Torricelli language of Sandaun Province, Papua New Guinea. Speaker counts differ by source. Wikipedia cites about 2,500 from the 2000 census. A Buffalo-led documentation project estimates about 1,000 speakers across 19 villages and hamlets. Verbix lists 2,134. Tok Pisin is the regional lingua franca. The search turned up no school materials in Yil.

\section{Primary Educational Resources}
No primers, textbooks or school programmes for Yil were found in search results.

\section{Linguistic and Reference Materials}
\begin{itemize}
\item Wikipedia, ``Yil language'': overview with a phonology sketch based on Martens and Tuominen. \url{https://en.wikipedia.org/wiki/Yil_language}
\item ELDP archive, ``Documentation and description of Yil, a Torricelli language of Papua New Guinea''. This is the first major documentation effort since the 1970s. Earlier sources were limited to Laycock's handwritten field notes and a short phonological sketch. \url{https://eldp.access.preservica.com/dk0700/}
\item Verbix endangered-languages page for Yil. \url{https://docs.verbix.com/EndangeredLanguages/Yil}
\end{itemize}

\section{Religious and Community Resources}
\begin{itemize}
\item Joshua Project profile of Yil in Papua New Guinea. It is a demographic profile only; no Yil scripture was confirmed. \url{https://joshuaproject.net/people_groups/16036}
\end{itemize}

\section{Digital and Audio Resources}
The ELDP project holds recordings (annotations cover about half of the elicited recordings), but access conditions were not verified here. See the ELDP archive link above.

\section{Curriculum Development Context}
Yil is described as moribund or under-documented, and no school use is known. Curriculum work would need new materials built from the documentation corpus and the involvement of community members, who should be consulted directly. Very little published material exists.

\end{document}